\textbf{Phần 1. Trắc nghiệm lựa chọn}
\setcounterpageref{ex}{0}

\begin{ex} %[4D13V69]
Phân số nào sau đây có tử số bằng 5 và mẫu số bằng 8?
\choice
{$\dfrac{8}{5}$}
{$5{,}8$}
{\True $\dfrac{5}{8}$}
{$8{,}5$}
\loigiai{
Phân số đó là $\dfrac{5}{8}$.\\
\textbf{Đáp án của câu hỏi là C}
}
\end{ex}

\begin{ex} %[4D13V69]
Phân số $\dfrac{12}{16}$ rút gọn về phân số tối giản là
\choice
{$\dfrac{6}{8}$}
{\True $\dfrac{3}{4}$}
{$\dfrac{2}{3}$}
{$\dfrac{4}{5}$}
\loigiai{
$\dfrac{12 : 4}{16 : 4} = \dfrac{3}{4}$.\\
\textbf{Đáp án của câu hỏi là B}
}
\end{ex}

\begin{ex} %[4D13V69]
Phân số nào sau đây bằng phân số $\dfrac{2}{3}$?
\choice
{$\dfrac{4}{9}$}
{$\dfrac{6}{10}$}
{$\dfrac{3}{2}$}
{\True $\dfrac{8}{12}$}
\loigiai{
$\dfrac{2 \times 4}{3 \times 4} = \dfrac{8}{12}$.\\
\textbf{Đáp án của câu hỏi là D}
}
\end{ex}

\begin{ex} %[4D13V69]
Trong các phân số sau, phân số nào lớn hơn 1?
\choice
{\True $\dfrac{9}{7}$}
{$\dfrac{3}{4}$}
{$\dfrac{5}{5}$}
{$\dfrac{7}{9}$}
\loigiai{
$\dfrac{9}{7} > 1$ (tử số lớn hơn mẫu số).\\
\textbf{Đáp án của câu hỏi là A}
}
\end{ex}

\begin{ex} %[4D13V69]
Điền dấu thích hợp vào chỗ chấm: $\dfrac{4}{7} \dots \dfrac{6}{7}$.
\choice
{$>$}
{$=$}
{\True $<$}
{Không so sánh được}
\loigiai{
$4 < 6 \Rightarrow \dfrac{4}{7} < \dfrac{6}{7}$.\\
\textbf{Đáp án của câu hỏi là C}
}
\end{ex}

\begin{ex} %[4D13V69]
Mẫu số chung nhỏ nhất của hai phân số $\dfrac{3}{4}$ và $\dfrac{5}{6}$ là
\choice
{$24$}
{\True $12$}
{$18$}
{$36$}
\loigiai{
12 chia hết cho cả 4 và 6.\\
\textbf{Đáp án của câu hỏi là B}
}
\end{ex}

\begin{ex} %[4D13V69]
Sắp xếp các phân số $\dfrac{1}{2}, \dfrac{2}{3}, \dfrac{3}{4}$ theo thứ tự từ bé đến lớn:
\choice
{$\dfrac{3}{4} < \dfrac{2}{3} < \dfrac{1}{2}$}
{$\dfrac{2}{3} < \dfrac{1}{2} < \dfrac{3}{4}$}
{$\dfrac{1}{2} < \dfrac{3}{4} < \dfrac{2}{3}$}
{\True $\dfrac{1}{2} < \dfrac{2}{3} < \dfrac{3}{4}$}
\loigiai{
$\dfrac{6}{12} < \dfrac{8}{12} < \dfrac{9}{12} \Rightarrow \dfrac{1}{2} < \dfrac{2}{3} < \dfrac{3}{4}$.\\
\textbf{Đáp án của câu hỏi là D}
}
\end{ex}

\begin{ex} %[4D13V69]
Tìm $x$ biết: $\dfrac{x}{18} = \dfrac{2}{3}$.
\choice
{\True $12$}
{$10$}
{$14$}
{$16$}
\loigiai{
$x = 2 \times 6 = 12$.\\
\textbf{Đáp án của câu hỏi là A}
}
\end{ex}

\begin{ex} %[4D13V69]
Một lớp có 30 học sinh, trong đó có 18 bạn nữ. Phân số chỉ số bạn nam so với cả lớp rút gọn tối giản là
\choice
{$\dfrac{3}{5}$}
{$\dfrac{12}{30}$}
{\True $\dfrac{2}{5}$}
{$\dfrac{2}{3}$}
\loigiai{
Nam: $30 - 18 = 12$ bạn $\Rightarrow \dfrac{12}{30} = \dfrac{2}{5}$.\\
\textbf{Đáp án của câu hỏi là C}
}
\end{ex}

\begin{ex} %[4D13V69]
Phân số nào sau đây tối giản?
\choice
{$\dfrac{15}{25}$}
{\True $\dfrac{9}{14}$}
{$\dfrac{14}{21}$}
{$\dfrac{8}{12}$}
\loigiai{
$\dfrac{9}{14}$ là phân số tối giản.\\
\textbf{Đáp án của câu hỏi là B}
}
\end{ex}

\begin{ex} %[4D13V69]
Viết phân số $\dfrac{4}{5}$ thành phân số có mẫu số bằng 100:
\choice
{$\dfrac{40}{100}$}
{$\dfrac{85}{100}$}
{$\dfrac{75}{100}$}
{\True $\dfrac{80}{100}$}
\loigiai{
$\dfrac{4 \times 20}{5 \times 20} = \dfrac{80}{100}$.\\
\textbf{Đáp án của câu hỏi là D}
}
\end{ex}

\begin{ex} %[4D13V69]
Tìm phân số tối giản $\dfrac{a}{b}$ bằng $\dfrac{4}{7}$ biết mẫu số hơn tử số 15 đơn vị.
\choice
{\True $\dfrac{20}{35}$}
{$\dfrac{12}{27}$}
{$\dfrac{16}{31}$}
{$\dfrac{24}{39}$}
\loigiai{
Hiệu số phần: $7 - 4 = 3$ phần $= 15 \Rightarrow 1$ phần $= 5$.\\
Tử: $4 \times 5 = 20$, Mẫu: $7 \times 5 = 35 \Rightarrow \dfrac{20}{35}$.\\
\textbf{Đáp án của câu hỏi là A}
}
\end{ex}

\textbf{Phần 2. Trắc nghiệm đúng sai}
\setcounterpageref{ex}{0}

\begin{ex} %[4D13V69]
Xét tính đúng sai về phân số:
\choiceTF
{\True $\dfrac{3}{5} = \dfrac{6}{10} = \dfrac{9}{15}$.}
{\True Phân số $\dfrac{7}{8}$ bé hơn 1.}
{Phân số $\dfrac{5}{4}$ bé hơn phân số $\dfrac{4}{5}$.}
{\True Phân số $\dfrac{5}{4}$ lớn hơn phân số $\dfrac{4}{5}$.}
\loigiai{
\begin{itemize}
    \item Dãy phân số bằng nhau đúng.
    \item $7/8 < 1$: Đúng.
    \item $5/4 > 1 > 4/5$. Ý 3 Sai, Ý 4 Đúng.
\end{itemize}
\begin{itemize}
    \item a) \textbf{Đúng}.
    \item b) \textbf{Đúng}.
    \item c) \textbf{Sai}.
    \item d) \textbf{Đúng}.
\end{itemize}
}
\end{ex}

\begin{ex} %[4D13V69]
Cho 3 phân số: $\dfrac{3}{4}, \dfrac{5}{6}, \dfrac{7}{12}$. Xét tính đúng sai:
\choiceTF
{\True Mẫu số chung nhỏ nhất của 3 phân số là 12.}
{\True Quy đồng được $\dfrac{9}{12}, \dfrac{10}{12}, \dfrac{7}{12}$.}
{Phân số lớn nhất là $\dfrac{3}{4}$.}
{\True Phân số lớn nhất là $\dfrac{5}{6}$.}
\loigiai{
\begin{itemize}
    \item MSC là 12: Đúng.
    \item $9/12, 10/12, 7/12$: Đúng.
    \item Lớn nhất là $5/6 = 10/12 \neq 3/4$. Ý 3 Sai, Ý 4 Đúng.
\end{itemize}
\begin{itemize}
    \item a) \textbf{Đúng}.
    \item b) \textbf{Đúng}.
    \item c) \textbf{Sai}.
    \item d) \textbf{Đúng}.
\end{itemize}
}
\end{ex}

\textbf{Phần 3. Trả lời ngắn}
\setcounterpageref{ex}{0}

\begin{ex} %[4D13V69]
Rút gọn phân số $\dfrac{25}{35}$ về tối giản (dạng a/b).
\shortans{5/7}
\loigiai{
\textbf{Đáp số:} $5/7$\\
$5/7$.
}
\end{ex}

\begin{ex} %[4D13V69]
Tìm $x$ biết $\dfrac{x}{15} = \dfrac{3}{5}$.
\shortans{9}
\loigiai{
\textbf{Đáp số:} $9$\\
$x = 3 \times 3 = 9$.
}
\end{ex}

\begin{ex} %[4D13V69]
Quy đồng mẫu số hai phân số $\dfrac{2}{3}$ và $\dfrac{3}{5}$ với mẫu số chung 15. Tìm tổng hai tử số mới.
\shortans{19}
\loigiai{
\textbf{Đáp số:} $19$\\
$10 + 9 = 19$.
}
\end{ex}

\begin{ex} %[4D13V69]
Tìm phân số lớn nhất trong các phân số: $\dfrac{3}{5}, \dfrac{4}{5}, \dfrac{2}{5}$ (dạng a/b).
\shortans{4/5}
\loigiai{
\textbf{Đáp số:} $4/5$\\
$4/5$.
}
\end{ex}

\begin{ex} %[4D13V69]
Tìm phân số bé nhất trong các phân số: $\dfrac{5}{6}, \dfrac{5}{7}, \dfrac{5}{8}$ (dạng a/b).
\shortans{5/8}
\loigiai{
\textbf{Đáp số:} $5/8$\\
$5/8$.
}
\end{ex}

\begin{ex} %[4D13V69]
Tìm chữ số $x$ để $\dfrac{1}{2} < \dfrac{x}{8} < \dfrac{3}{4}$.
\shortans{5}
\loigiai{
\textbf{Đáp số:} $5$\\
$4/8 < x/8 < 6/8 \Rightarrow x = 5$.
}
\end{ex}
